% 第 15 章 SMVA
\bichapter{SMVA 基于图的同时多电压分析}{SMVA Graph-Based Simultaneous Multivoltage Analysis}
\label{chap:smva}

基于图的 SMVA（simultaneous multivoltage analysis）在单次分析运行中，同时考虑全设计所有路径的各路径供电电压的所有组合。关于 SMVA，参见：
\begin{itemize}
  \item SMVA 概述
  \item 配置 SMVA 分析
  \item 在 SMVA 分析中报告时序路径
  \item DVFS 场景
  \item 将 SMVA 与其他 PrimeTime 特性配合使用
  \item 使用示例
  \item 支持 DVFS 场景的命令与属性
  \item 特性兼容性
\end{itemize}

% ============================================================
\section{SMVA 概述}
\subsection*{Overview of SMVA}

多电压设计中的时序路径可能跨越电源域。当各电源域可在多种供电电压下工作时，传统分析无法准确反映各种电压组合对每条路径的影响：每个域只使用一个电压无法涵盖所有组合；对每个域使用 min/max（OCV）电压又会过于悲观，因为同一域内的不同单元不可能同时工作在不同电压下。

基于图的同时多电压分析（SMVA）消除该悲观度：域内路径在该域所有电压级分析，跨域路径分析所穿越域的所有电压组合，单次运行完成。

SMVA 对采用动态电压与频率缩放（DVFS）技术的设计尤其有用——芯片运行中根据任务需求调整供电电压与频率，高负载时更高功耗，其他时候更低功耗。提供 scaling 库组并指定各电源域供电电压即可获得准确时序结果。

\subsection{SMVA 要求}
\subsubsection*{SMVA Requirements}

基于图 SMVA 分析要求：
\begin{itemize}
  \item 设计须有多电压功耗意图信息（见第~\ref{chap:mv}~章「多电压分析要求」）
  \item SMVA 特性由 PrimeTime-ADV-PLUS 许可证启用；所需许可证数量取决于分析中使用的电压级数与 DVFS 场景数（见「SMVA 许可证」）
\end{itemize}

\subsection{SMVA 电源域分析}
\subsubsection*{SMVA Power Domain Analysis}

启用 SMVA 后，基于图分析在单次运行中同时考虑每条路径的所有供电电压组合（跨域与域内），排除各路径与路径段的不可能组合。

\figplaceholder{Figure 143: Supply Combinations for Cross-Domain and Within-Domain Paths}{跨域与域内路径的供电组合}{fig:smva-supply-combo}

域 A 与 B 的供电可独立设为低或高电压。从 FF2（蓝）与 FF5（紫）开始的域内路径分别在各自域的低、高供电电压下分析。从 FF1（红）开始的跨域路径在两域供电的四种可能组合下分析。

给定供电条件下的每条路径都作为单独路径进行收集和报告；因此，同一条物理路径可由 \cmd{report\_timing -nworst} 针对不同供电组合分别报告。数据路径或时钟路径中的任一 timing arc 跨越电源域边界时，该路径即为跨域路径。UPF 基础设施确定电源域归属；每个单元的 PG 引脚都应连接到供电网，每个端口也应定义供电网。工具既识别 Verilog PG 网表中的连接，也识别 UPF \cmd{connect\_supply\_net} 指定的连接，其中 Verilog PG 网表是首选方式。若路径中有对象缺少供电网信息，工具会根据设计和库信息推断最可能的供电，并使用该供电的单一电压值；可用 \cmd{check\_timing -include unconnected\_pg\_pins} 和 \cmd{check\_timing -include supply\_net\_voltage} 报告缺失信息。

% ============================================================
\section{配置 SMVA 分析}
\subsection*{Configuring SMVA Analysis}

在首次时序更新前：
\begin{enumerate}
  \item 确保多电压功耗意图与 scaling 库组已配置（第~\ref{chap:mv}~章）
  \item 启用 SMVA：\cmd{set\_app\_var timing\_enable\_cross\_voltage\_domain\_analysis true}
  \item 用 \cmd{get\_supply\_groups} 获取供电组，并用 \cmd{set\_voltage\_levels} 定义命名电压参考
  \item 用 \cmd{set\_voltage} 为各供电组的命名电压参考设置电压值
\end{enumerate}

可选地，将 \cmd{timing\_cross\_voltage\_domain\_analysis\_mode} 设为 \texttt{full}（默认）、\texttt{only\_crossing}、\texttt{exclude\_crossing} 或 \texttt{legacy}，以控制计算 slack 以及报告/ECO 命令所考虑的路径范围。

\subsection{定义命名 SMVA 电压参考}
\subsubsection*{Defining Named SMVA Voltage References}

用 \cmd{get\_supply\_groups} 获取具有指定供电连接的供电组集合；用 \cmd{set\_voltage\_levels -reference\_names} 为供电组定义命名电压参考；再用 \cmd{set\_voltage -reference\_name} 将具体电压值与各参考名称关联。

例如，顶层含下级块 B1、B2，各块供电组均可在 0.8 V 与 1.0 V 下工作：
\begin{lstlisting}
set_app_var timing_enable_cross_voltage_domain_analysis true
set sg_b1 [get_supply_groups B1/VDD]
set sg_b2 [get_supply_groups B2/VDD]
set_voltage_levels -reference_names {lo hi} $sg_b1
set_voltage_levels -reference_names {lo hi} $sg_b2
set_voltage -reference_name lo -object_list $sg_b1 0.8
set_voltage -reference_name hi -object_list $sg_b1 1.0
set_voltage -reference_name lo -object_list $sg_b2 0.8
set_voltage -reference_name hi -object_list $sg_b2 1.0
update_timing
\end{lstlisting}

上述 SMVA 专用命令与选项仅在启用 \cmd{timing\_enable\_cross\_voltage\_domain\_analysis} 后可用。

\begin{noteBox}
SMVA 配置须在首次 \cmd{update\_timing} 之前完成。
\end{noteBox}

% ============================================================
\section{在 SMVA 分析中报告时序路径}
\subsection*{Reporting Timing Paths in an SMVA Analysis}

SMVA 分析中，\cmd{report\_timing} 在每条路径报告中增加一行，显示该路径的供电条件（各供电网组上的命名电压参考）。

\figplaceholder{Figure 144: Reference Voltage Conditions}{参考电压条件}{fig:smva-ref-voltage}

可用 \opt{-supply\_net\_group} 查看路径上各点的供电组信息。\cmd{get\_timing\_paths} 返回的 \texttt{timing\_path} 集合对象含 \texttt{dvfs\_scenario} 属性，其 \texttt{description} 子属性描述路径的 DVFS 场景。

如「SMVA 电源域分析」所述，给定电压条件下的每条路径作为单独路径。默认 SMVA 同时考虑域内与跨域路径；\cmd{report\_timing} 和 \cmd{get\_timing\_paths} 的 \opt{-domain\_crossing} 可取 \texttt{all}（默认）、\texttt{only\_crossing} 或 \texttt{exclude\_crossing}，分别选择所有路径、仅跨域路径或仅域内路径。

% ============================================================
\section{DVFS 场景}
\subsection*{DVFS Scenarios}

对使用 DVFS 的设计，可用 DVFS 场景功能在 SMVA 之上进一步指定频率缩放与其他工作条件。

\begin{noteBox}
DVFS 场景是对 SMVA 的扩展；单独 SMVA 仅变化供电电压。
\end{noteBox}

\subsection{DVFS 场景概念}
\subsubsection*{DVFS Scenario Concepts}

SMVA 允许为各供电网组定义多种供电电压。DVFS 是调整供电电压与频率以降低平均功耗的设计技术。SMVA 的 DVFS 场景特性允许在电压组合之外指定频率与其他条件，以正确分析这些设计。

\subsection{DVFS 场景集合对象}
\subsubsection*{DVFS Scenario Collection Objects}

DVFS 场景表示为 \texttt{dvfs\_scenario} 集合对象，描述各供电网组上的命名电压参考组合。例如 \texttt{\{VDD1:lo VDDT:lo\}} 表示 VDD1 与 VDDT 均在低电压参考。

\cmd{get\_dvfs\_scenarios} 通过一组或多组 \opt{-supply\_group}/\opt{-reference\_names} 选项对构造 \texttt{dvfs\_scenario} 集合对象。供电组与命名电压参考的顺序不重要；对象按需构造，工具不会在内存中预先枚举所有局部和完整条件组合。

\begin{noteBox}
\texttt{dvfs\_scenario} 集合对象在查询时构造与使用；文档中简称为「DVFS 场景」。
\end{noteBox}

\subsection{传播的 DVFS 场景}
\subsubsection*{Propagated DVFS Scenarios}

每条时序路径都有一组完整描述其可能电压条件的 DVFS 场景；频率等条件则通过应用于这些场景的时序约束体现。

\figplaceholder{Figure 145: Fully Describing DVFS Scenarios for Individual Timing Paths}{完全描述单条时序路径的 DVFS 场景}{fig:smva-dvfs-path}

\figplaceholder{Figure 146: The Propagated DVFS Scenarios for a Design}{设计的传播 DVFS 场景}{fig:smva-dvfs-design}

获取设计中每条路径的完全描述场景即得到传播 DVFS 场景集合。例如三供电网组（VDDT、VDD1、VDD2）各两电压级，传播场景为所有合法组合。

\subsection{用 DVFS 场景控制命令作用域}
\subsubsection*{Using DVFS Scenarios to Control Command Scope}

对约束或报告命令使用 DVFS 场景时，命令仅作用于与该场景兼容的路径。

\figplaceholder{Figure 147: Applying a DVFS Scenario to a Command}{将 DVFS 场景应用于命令}{fig:smva-dvfs-cmd}

示例：对 \cmd{report\_timing} 指定 \texttt{\{VDD1:lo VDDT:lo\}} 会排除 VDD1 或 VDDT 使用其他命名电压参考的不兼容路径；但完全位于 VDD2 内的路径不会被排除，因为该命令场景未对 VDD2 施加条件。也就是说，命令会排除不兼容场景，并非只纳入字面列出的供电组。

\subsection{查询 DVFS 场景}
\subsubsection*{Querying DVFS Scenarios}

DVFS 场景以集合对象引用。

\subsubsection{为供电网组条件构造 DVFS 场景}
\paragraph*{Constructing DVFS Scenarios for Supply Group Conditions}

指定每个供电网组上的一个命名电压参考，返回具有指定条件的 DVFS 场景集合，用于传递给约束或报告命令。

\subsubsection{获取所有传播的 DVFS 场景}
\paragraph*{Obtaining All Propagated DVFS Scenarios}

\cmd{get\_dvfs\_scenarios -propagated} 返回当前设计的完整传播 DVFS 场景集合，可用于设计探索；若设计尚未进行逻辑更新，该命令会触发逻辑设计更新。并非所有支持 \opt{-dvfs\_scenarios} 的命令都接受包含多个场景对象的集合。

\subsection{将 DVFS 场景应用于命令与属性}
\subsubsection*{Applying DVFS Scenarios to Commands and Attributes}

默认情况下，支持 DVFS 场景的命令与属性考虑所有传播 DVFS 场景。

\subsubsection{应用于单个命令}
\paragraph*{Applying DVFS Scenarios to Individual Commands}

带 \opt{-dvfs\_scenarios} 的命令允许显式指定 DVFS 场景；显式指定优先于脚本或上下文级设置。

\subsubsection{应用于属性查询}
\paragraph*{Applying DVFS Scenarios to Attribute Queries}

查询特定 DVFS 场景下的属性值时，应在属性名后的圆括号中放入 \texttt{dvfs\_scenario} 集合对象，例如 \cmd{get\_attribute \$pin pin\_capacitance\_max(\$low\_low)}。若集合含多个场景，工具按属性类型返回最坏情况值；若字符串等值无法合并，则发出 ATTR-11 警告。不支持 DVFS 场景的属性查询会报 ATTR-3 错误。属性下标优先于脚本或上下文级 DVFS 场景设置。

\subsubsection{应用于脚本}
\paragraph*{Applying DVFS Scenarios to Scripts}

用 \cmd{read\_sdc -dvfs\_scenarios} 或 \cmd{source -dvfs\_scenarios} 对整个 SDC 或脚本应用 DVFS 场景，等效于对每个命令使用 \opt{-dvfs\_scenarios}。

\begin{noteBox}
并非所有 SDC 命令都支持 \opt{-dvfs\_scenarios}。为避免含义不明确，带 DVFS 场景的约束脚本中应只使用明确支持该选项的命令。
\end{noteBox}

\subsubsection{为命令执行设置 DVFS 场景上下文}
\paragraph*{Setting a DVFS Scenario Context for Command Execution}

\cmd{push\_dvfs\_scenario -dvfs\_scenario} 为后续支持 DVFS 的命令和属性查询建立场景上下文，\cmd{pop\_dvfs\_scenario} 移除最近压入的上下文。上下文可以嵌套，但不支持在 DVFS 上下文中执行时序更新。


% ============================================================
\section{将 SMVA 与其他 PrimeTime 特性配合使用}
\subsection*{Using SMVA Analysis With Other PrimeTime Features}

\subsection{与 OCV/AOCV/POCV 配合}
\subsubsection*{OCV, AOCV, and POCV}

SMVA 处理供电电压组合，而 OCV、AOCV、POCV 处理片上变异，两者描述的是不同分析维度，可同时启用。

\subsection{SMVA 流程中的电压缩放时序降额}
\subsubsection*{Voltage-Scaled Timing Derates in SMVA Flows}

\cmd{set\_timing\_derate} 可在不同对象作用域上指定 early/late 时序调整。在 SMVA 流程中，可应用与电压相关的时序降额，并由工具在分析中按需动态缩放。

要启用该特性，设置：
\begin{lstlisting}
pt_shell> set_app_var \
  timing_enable_derate_scaling_for_library_cells_compatibility false
pt_shell> set_app_var \
  timing_enable_derate_scaling_interpolation_for_library_cells true
\end{lstlisting}

启用后，若在 scaling 库组的库单元上设置时序降额，工具会按各 DVFS 场景中计算得到的实例电压缩放降额。例如：
\begin{lstlisting}
define_scaling_lib_group {lib_1.1V.db lib_1.0V.db lib_0.9V.db}
set_timing_derate -late [get_lib_cell lib_1.1V/INV1] 1.07
set_timing_derate -late [get_lib_cell lib_1.0V/INV1] 1.10
set_timing_derate -late [get_lib_cell lib_0.9V/INV1] 1.15
set_voltage_levels sg1 -reference_names {hi lo}
set_voltage -object_list sg1 -reference_name lo 0.90  ;# exact derate
set_voltage -object_list sg1 -reference_name hi 1.05  ;# scaled derate
\end{lstlisting}

在该示例中，对 INV1 单元：
\begin{itemize}
  \item \texttt{sg1:lo} — 0.90~V 与 0.9~V 库精确匹配，应用该库上的降额 1.15
  \item \texttt{sg1:hi} — 1.05~V 与任一库均非精确匹配，取最邻近两库（1.0~V 与 1.1~V）上降额 1.07 与 1.10 中的更差者 1.10
\end{itemize}

\cmd{set\_timing\_derate} 还提供 \opt{-dvfs\_scenarios}，可直接指定降额所适用的 DVFS 场景；仅支持叶级单元实例，且场景只能引用该单元相关供电。与常规降额相同，作用于单元实例的设置会覆盖对应库单元上的更一般设置。

\subsection{与多场景分析（DMSA）配合}
\subsubsection*{Distributed Multi-Scenario Analysis}

DVFS 场景与 DMSA 场景不是同一概念：DVFS 电压/频率场景在一次 SMVA 运行中同时分析，而 DMSA 的工作条件与约束场景在分布式运行中分析。

\subsection{与电平转换器、隔离与保持}
\subsubsection*{Level Shifters, Isolation, and Retention}

SMVA 尊重 UPF 指定的电平转换器、隔离单元与保持寄存器约束。跨域路径经过电平转换器时，在相应电压组合下分析转换器延时。

\subsection{与 IR Drop 标注配合}
\subsubsection*{IR Drop Annotation}

SMVA 流程可以使用实例级 \cmd{set\_voltage -cell -pg\_pin\_name} 电压标注。每个 DVFS 场景按其供电条件和实例级标注计算，不应把实例级标注简单理解为对命名参考电压作数值相加。

\subsection{与信号完整性分析配合}
\subsubsection*{Signal Integrity Analysis}

PrimeTime SI 可在 SMVA 电压条件下考虑 aggressor 与 victim 的相应电源轨电压。噪声分析有明确限制：\cmd{update\_noise -dvfs\_scenarios} 一次只接受一个完整描述设计供电条件的 DVFS 场景，\cmd{report\_noise} 报告最近一次更新的该场景结果；分析多个场景时须逐个执行。

\subsection{与路径分析配合}
\subsubsection*{Path-Based Analysis}

\cmd{report\_timing} 与 \cmd{get\_timing\_paths} 可用 \opt{-dvfs\_scenarios} 将报告范围限制到兼容的 DVFS 场景，并可结合 \opt{-domain\_crossing} 选择跨域或域内路径。

% ============================================================
\section{使用示例}
\subsection*{Usage Examples}

\subsection{两域双电压 SMVA}
\subsubsection*{Two Domains, Two Voltages Each}

\begin{lstlisting}
# Define supply groups and voltage references
set_app_var timing_enable_cross_voltage_domain_analysis true
set sg_core [get_supply_groups VDD_CORE]
set sg_io [get_supply_groups VDD_IO]
set_voltage_levels -reference_names {v0p9 v1p0} $sg_core
set_voltage_levels -reference_names {v0p9 v1p0} $sg_io
set_voltage -reference_name v0p9 -object_list $sg_core 0.9
set_voltage -reference_name v1p0 -object_list $sg_core 1.0
set_voltage -reference_name v0p9 -object_list $sg_io 0.9
set_voltage -reference_name v1p0 -object_list $sg_io 1.0
update_timing
report_timing -max_paths 10
\end{lstlisting}

\subsection{带 DVFS 场景的报告}
\subsubsection*{Reporting With DVFS Scenarios}

\begin{lstlisting}
set scen [get_dvfs_scenarios \
          -supply_group $sg_core -reference_names v0p9 \
          -supply_group $sg_io -reference_names v1p0]
report_timing -dvfs_scenarios $scen -max_paths 5
get_attribute [get_pins U1/ZN] pin_capacitance_max($scen)
\end{lstlisting}

\subsection{跨域路径过滤}
\subsubsection*{Cross-Domain Path Filtering}

\begin{lstlisting}
report_timing -domain_crossing only_crossing -max_paths 20
report_timing -domain_crossing exclude_crossing -max_paths 20
\end{lstlisting}

% ============================================================
\section{支持 DVFS 场景的命令与属性}
\subsection*{Commands and Attributes With DVFS Scenario Support}

支持 \opt{-dvfs\_scenarios} 或 DVFS 场景下标的命令包括（非穷尽）：
\begin{itemize}
  \item \cmd{report\_timing}、\cmd{report\_constraint}、\cmd{get\_timing\_paths}
  \item \cmd{set\_false\_path}、\cmd{set\_multicycle\_path}、\cmd{set\_max\_delay}、\cmd{set\_min\_delay}
  \item \cmd{set\_clock\_uncertainty}、\cmd{set\_input\_delay}、\cmd{set\_output\_delay}
  \item 多种 pin、port、net、clock 与 \texttt{pg\_pin\_info} 属性
\end{itemize}

\begin{table}[htbp]
\centering
\caption{支持 DVFS 场景的属性（节选）}
\label{tab:smva-attrs}
\small
\begin{tabular}{|l|l|l|}
\hline
\tblhead{属性名} & \tblhead{类} & \tblhead{说明} \\
\hline
\texttt{dvfs\_scenario} & \texttt{timing\_path} & 路径的 DVFS 场景 \\
\texttt{pin\_capacitance\_max}, \texttt{pin\_capacitance\_min} & \texttt{pin}, \texttt{port}, \texttt{net} & 可按 DVFS 场景查询 \\
\texttt{power\_rail\_voltage\_max}, \texttt{power\_rail\_voltage\_min} & \texttt{pin}, \texttt{port} & 可按 DVFS 场景查询 \\
\hline
\end{tabular}
\end{table}

% ============================================================
\section{特性兼容性}
\subsection*{Feature Compatibility}

\begin{table}[htbp]
\centering
\caption{SMVA 特性限制}
\label{tab:smva-compat}
\small
\begin{tabularx}{\textwidth}{|l|X|}
\hline
\tblhead{特性} & \tblhead{限制} \\
\hline
噪声分析 & 不支持一次分析多个 DVFS 场景，必须逐场景更新 \\
ETM 生成 & 不支持生成单个汇总所有 SMVA 条件的 ETM；工具按相关活动场景分别生成模型 \\
Reduced-resource ECO & 不兼容 \\
SDF 回标 & 不支持为不同 DVFS 场景使用不同的 SDF 延时数据 \\
PLL generated clocks & 不支持 \cmd{create\_generated\_clock} 的 \opt{-pll\_feedback} 与 \opt{-pll\_output} \\
高级属性用法 & 某些用法不兼容，须以属性是否明确支持 DVFS 场景为准 \\
\hline
\end{tabularx}
\end{table}

\begin{noteBox}
SMVA 会增加分析状态空间以及运行时间和内存消耗。电压级、供电组与 DVFS 场景较多时，应评估状态组合及许可证需求。
\end{noteBox}

\begin{seeAlsoBox}
第~\ref{chap:mv}~章多电压流程；第~\ref{chap:var}~章片上变异；PrimeTime 许可证章节中的 SMVA 许可证说明。
\end{seeAlsoBox}
